\chapter{Multiple-Qubit Systems and Cryptography --- Solutions}

\noindent\textbf{SP6:} Reading---Rieffel--Polak \S3; Justin Ch.~13 (Cryptography). \textbf{Quiz:} Rieffel 3.1--3.6 (for ``all $n>1$'' arguments, only $n=2$ is required); Justin \S13.1--5. \textbf{Exam:} Same Rieffel; Justin---complement of quiz (from \S13.6 on). \textit{Octave:} \texttt{ex11\_sp6.m}.

%% -------------------------------------------------------------------------
\section*{Rieffel--Polak: Exercises 3.1--3.6}

\begin{exercise}{3.1}
Let $V$ be a vector space with basis $\{(1,0,0),(0,1,0),(0,0,1)\}$. Give two different bases for $V\otimes V$.
\end{exercise}
\begin{solution}
Identify $V=\C^3$ with standard orthonormal basis $\ket{e_1},\ket{e_2},\ket{e_3}$.

\textbf{Basis A (tensor products of standard basis vectors):}
\[
\mathcal{B}_A=\{\ket{e_i}\otimes\ket{e_j}: i,j\in\{1,2,3\}\},
\]
nine vectors, orthonormal for the product inner product.

\textbf{Basis B (tensor products with a Fourier basis on the first factor):} Let $\omega=e^{2\pi i/3}$ and define
\[
\ket{f_k}=\tfrac{1}{\sqrt{3}}\sum_{j=0}^{2}\omega^{jk}\ket{e_{j+1}},\qquad k=0,1,2.
\]
Then $\{\ket{f_0},\ket{f_1},\ket{f_2}\}$ is an orthonormal basis of $V$ (unitary DFT matrix). Set
\[
\mathcal{B}_B=\{\ket{f_k}\otimes\ket{e_\ell}: k,\ell\in\{0,1,2\}\}
\]
(indexing $e_j$ with $j=1,2,3$ corresponding to $\ell$). This is again an orthonormal basis of $V\otimes V$, and $\mathcal{B}_B\neq \mathcal{B}_A$ because $\ket{f_0}=\tfrac{1}{\sqrt{3}}(\ket{e_1}+\ket{e_2}+\ket{e_3})$ is not a basis vector of $\mathcal{B}_A=\{\ket{e_i}\otimes\ket{e_j}\}$.

\paragraph{Remark.} Any pair of orthonormal bases $\{u_i\}$, $\{v_j\}$ of $V$ yields an orthonormal basis $\{u_i\otimes v_j\}$ of $V\otimes V$; choosing different bases gives different tensor-product bases.
\end{solution}

\begin{exercise}{3.2}
Show by example that a linear combination of entangled states is not necessarily entangled.
\end{exercise}
\begin{solution}
Let $\ket{\Phi^+}=\tfrac{1}{\sqrt{2}}(\ket{00}+\ket{11})$ and $\ket{\Phi^-}=\tfrac{1}{\sqrt{2}}(\ket{00}-\ket{11})$. Both are entangled (Schmidt rank $2$). Their sum is
\[
\ket{\Phi^+}+\ket{\Phi^-}=\tfrac{2}{\sqrt{2}}\ket{00}=\sqrt{2}\,\ket{0}\otimes\ket{0},
\]
which is a product state (up to normalization). Thus a linear combination of entangled pure states can be separable.
\end{solution}

\begin{exercise}{3.3}
Show that $\ket{W_n}$ is entangled for every $n>1$ (course: prove $n=2$).
\end{exercise}
\begin{solution}
For $n=2$, $\ket{W_2}=\tfrac{1}{\sqrt{2}}(\ket{01}+\ket{10})$. Suppose $\ket{W_2}=\ket{a}\otimes\ket{b}$ with $\ket{a}=a_0\ket{0}+a_1\ket{1}$, $\ket{b}=b_0\ket{0}+b_1\ket{1}$. Expanding,
\[
\ket{a}\otimes\ket{b}=a_0b_0\ket{00}+a_0b_1\ket{01}+a_1b_0\ket{10}+a_1b_1\ket{11}.
\]
Matching $\ket{W_2}$: coefficients of $\ket{00}$ and $\ket{11}$ are $0$, so $a_0b_0=0$ and $a_1b_1=0$. If $a_0=0$ then $a_1\neq 0$ (unit vector), so $b_1=0$, hence $b_0=1$; then the $\ket{01}$ coefficient is $0$, contradiction. If $b_0=0$ then $b_1\neq 0$, so $a_1=0$, hence $a_0=1$; then the $\ket{10}$ coefficient is $0$, contradiction. Hence no product decomposition; $\ket{W_2}$ is entangled.
\end{solution}

\begin{exercise}{3.4}
Show that $\ket{\mathrm{GHZ}_n}$ is entangled for every $n>1$ (course: prove $n=2$).
\end{exercise}
\begin{solution}
For $n=2$, $\ket{\mathrm{GHZ}_2}=\tfrac{1}{\sqrt{2}}(\ket{00}+\ket{11})=\ket{\Phi^+}$, the Bell state. If $\ket{\Phi^+}=\ket{a}\otimes\ket{b}$, comparing coefficients of $\ket{01}$ and $\ket{10}$ gives $a_0b_1=a_1b_0=0$. If $a_0=0$ then $a_1\neq 0$ forces $b_0=0$ and $b_1=1$, giving coefficient $0$ on $\ket{00}$, contradiction. If $b_0=0$ similarly $a_1=0$ and $\ket{10}$ vanishes. Hence entangled.
\end{solution}

\begin{exercise}{3.5}
Is $\ket{\psi}=\tfrac{1}{\sqrt{2}}(\ket{0}\ket{+}+\ket{1}\ket{-})$ entangled?
\end{exercise}
\begin{solution}
Expand in the computational basis $\{\ket{00},\ket{01},\ket{10},\ket{11}\}$:
\[
\ket{+}=\tfrac{1}{\sqrt{2}}(\ket{0}+\ket{1}),\quad \ket{-}=\tfrac{1}{\sqrt{2}}(\ket{0}-\ket{1}),
\]
so $\ket{0}\ket{+}=\tfrac{1}{\sqrt{2}}(\ket{00}+\ket{01})$ and $\ket{1}\ket{-}=\tfrac{1}{\sqrt{2}}(\ket{10}-\ket{11})$. Therefore
\[
\ket{\psi}=\tfrac{1}{2}(\ket{00}+\ket{01}+\ket{10}-\ket{11}).
\]
If $\ket{\psi}=\ket{a}\otimes\ket{b}$ with $\ket{a}=a_0\ket{0}+a_1\ket{1}$, $\ket{b}=b_0\ket{0}+b_1\ket{1}$, then the coefficients are $a_0b_0=1/2$, $a_0b_1=1/2$, $a_1b_0=1/2$, $a_1b_1=-1/2$. From $a_0b_0=a_0b_1$ with $a_0\neq 0$ we get $b_0=b_1$; from $a_1b_0=-a_1b_1$ with $a_1\neq 0$ we get $b_0=-b_1$, impossible unless $b_0=b_1=0$, contradicting $|b_0|^2+|b_1|^2=1$. Hence \textbf{yes, entangled}.
\end{solution}

\begin{exercise}{3.6}
If someone asks whether $\ket{+}$ is entangled, what will you say?
\end{exercise}
\begin{solution}
\textbf{Entanglement} is a property of \emph{composite} systems (multi-partite states in a tensor product). The state $\ket{+}$ lives in a \emph{single} qubit space $\C^2$. It is \emph{not} a superposition in the $\{\ket{0},\ket{1}\}$ basis, but the question ``is $\ket{+}$ entangled?'' is \textbf{not meaningful} unless a tensor-product decomposition into subsystems is specified. Conclusion: \textbf{``No---$\ket{+}$ is a single-qubit state; entanglement is not defined for that setting.''}
\end{solution}

%% -------------------------------------------------------------------------
\section*{Justin Wyss-Gallifent: Quiz (\S13.1--5)}

\begin{exercise}{13.1}
Encrypt the stream \texttt{10110100010100101} using the key \texttt{10111}.
\end{exercise}
\begin{solution}
Bit-by-bit addition mod $2$ (XOR) with the key repeated periodically. Index $i$ uses key digit $k_i = k_{((i-1)\bmod 5)+1}$.

Plaintext: \texttt{10110100010100101}\\
Key (rep): \texttt{10111101111011110}\\
Ciphertext: \texttt{00001001101111011}

\paragraph{Check (first five bits):} $1\!+\!1\equiv 0$, $0\!+\!0\equiv 0$, $1\!+\!1\equiv 0$, $1\!+\!1\equiv 0$, $0\!+\!1\equiv 1$. \textit{Octave:} \texttt{ex11\_sp6.m}.
\end{solution}

\begin{exercise}{13.2}
Encrypt the stream \texttt{110110100100000101100101} using the key \texttt{110101}.
\end{exercise}
\begin{solution}
Key length $6$. Ciphertext (from \texttt{ex11\_sp6.m}):
\[
\texttt{000011010001110000010000}.
\]
\end{solution}

\begin{exercise}{13.3}
First $30$ digits of the key defined by $\{[1;1;0;0],[1;0;0;1]\}$.
\end{exercise}
\begin{solution}
Here $i=4$, $\bar{s}=[1;1;0;0]$, $\bar{c}=[1;0;0;1]$. Justin's recurrence gives
\[
x_n \equiv x_{n-4} + 0\cdot x_{n-3} + 0\cdot x_{n-2} + 1\cdot x_{n-1} \pmod 2,
\]
i.e.\ $x_n\equiv x_{n-4}+x_{n-1}\pmod 2$. The first $30$ bits are
\[
\texttt{110010001111010110010001111010}.
\]
The period (first $200$ bits) is $15$. \textit{Octave:} \texttt{ex11\_sp6.m}.
\end{solution}

\begin{exercise}{13.4}
First $30$ digits for $\{[1;0;1;1],[1;0;0;1]\}$.
\end{exercise}
\begin{solution}
Same recurrence $x_n\equiv x_{n-4}+x_{n-1}\pmod 2$ with initial $\bar{s}=[1;0;1;1]$. First $30$ bits:
\[
\texttt{101100100011110101100100011110}.
\]
Period $15$ (same recurrence, different phase). \textit{Octave:} \texttt{ex11\_sp6.m}.
\end{solution}

\begin{exercise}{13.5}
Brute force: recursion for fragment \texttt{1011100101}; find period.
\end{exercise}
\begin{solution}
Try smallest $i\ge 2$. For $i=3$, the recurrence is $x_n\equiv x_{n-3}+x_{n-2}\pmod 2$ (coefficients $\bar{c}=[1;1;0]$). This is consistent with the entire fragment $x_1\ldots x_{10}$.

\textbf{Finding the period:} Generate the key from any compatible initial triple (e.g.\ the first three bits of the fragment) and $\bar{c}$; the period is the smallest $T$ with $x_{k+T}=x_k$ for all $k$ in the generated sequence. (See Justin \S13.6 for the general algorithm; \texttt{ex11\_sp6.m} confirms the shortest recursion length $i=3$ for this fragment.)
\end{solution}

%% -------------------------------------------------------------------------
\section*{Justin Wyss-Gallifent: Exam (complement of quiz)}

\begin{exercise}{13.6}
Brute force for fragment \texttt{0011101001}.
\end{exercise}
\begin{solution}
Shortest linear recursion has length $i=3$ with $\bar{c}=[1;0;1]$ (Octave: \texttt{ex11\_sp6.m}):
\[
x_n \equiv x_{n-3} + x_{n-1} \pmod 2.
\]
\end{solution}

\begin{exercise}{13.7--13.12}
(Refined brute force on longer fragments; ciphertext attacks in 13.13--14.)
\end{exercise}
\begin{solution}
\textbf{Method:} Follow Justin \S13.6.3 (determinant test for shortest length) and \S13.7 (solve for $\bar{c}$ mod $2$). For each fragment, build the matrix equations $x_{n}=\sum_{j}c_j x_{n-j}$ and verify consistency. For 13.13--14, obtain key bits from ciphertext$\oplus$plaintext, then run the same search. \textit{Implementation:} extend \texttt{ex11\_sp6.m} with the same brute-force loop over larger $i$ and fragment lengths; for large fragments prefer Justin's refined search to reduce cases.
\end{solution}

\begin{exercise}{13.15--13.20}
(Pseudorandom digits, determinants, short recursions.)
\end{exercise}
\begin{solution}
Use the recursive-key generator and \S13.6--13.7: express unknowns as solutions of $A\bar{c}\equiv\bar{b}\pmod 2$; for 13.18, interpret the determinant pattern as in \S13.6.3 (shortest length guess). For 13.20, exhibit a length-$2$ recurrence that generates the same infinite periodic key as a length-$4$ specification (reduce the linear system).
\end{solution}
